\bookchapter{第十六章：賬號體系與雲端之務}

\chapterintro{吾有一友，去年手機被竊。其第一時間所思非“手機值幾何”，乃“吾內之照片還能覓回否”。甚憾，其未開雲備份。手機乃身外之物，然賬號內之數據，乃真失即再無回者。此章即論何以給君之數字生活上一真正之保險。}

\booksection{一、君之手機賬號，究竟管幾何物？}

衆人覺賬號乃用以下載App者。實則，一手機賬號背後所牽之物，遠較君所想爲多：

\begin{itemize}
  \item \textbf{通訊錄}：君存了幾年甚至十幾年之電話號碼。
  \item \textbf{照片與視頻}：自孩子出生至畢業典禮，自第一次約會至結婚旅行。
  \item \textbf{備忘錄}：君隨手記之靈感、日記、賬號密碼（對，多人將密碼記於備忘錄中）。
  \item \textbf{App購買記錄}：花了數百塊買之付費App與內購。
  \item \textbf{查找手機}：手機失之後，此乃唯一能助君定位、鎖機、遠程抹數據之工具。
\end{itemize}

若君用蘋果手機，以上所有物皆綁於君之Apple ID上。若君用安卓，即綁於君之華爲/小米/OPPO/vivo賬號上，尚有君之谷歌賬號。此等賬號即君數字生活之“戶口本”，失戶口本，諸多物即說不清矣。

\booksection{二、三必爲之事，今即爲之}

\begin{enumerate}
  \item \textbf{啓雙重認證}：無論君用蘋果抑或安卓，第一事即往賬號設置中將"雙重認證"或"兩步驗證"開之。此物如爲君之門加一反鎖——光有密碼不夠，尚需君之手機收一驗證碼。縱有人於網上猜到君之密碼，無君之手機，彼亦進不去。
  \item \textbf{綁一可靠之恢復方式}：雙重認證甚安全，然亦有一險——君己之手機失奈何？故系統將提示君設"恢復密鑰"或"備用郵箱/手機號"。此乃君於緊急之時取回賬號之唯一途徑，最好記於紙上，勿存手機中。
  \item \textbf{確認"查找手機"已啓}：此能於手機失之後能助君定位、鎖定、播放聲音、甚至遠程擦除所有數據。乃最後一道防線。多人買手機之時跳過此步，後欲用時覺未開，即晚矣。
\end{enumerate}

\booksection{三、iCloud：蘋果用戶之“空氣”}

用蘋果手機之人，遲早須對一靈魂之問：iCloud空間滿奈何？

用蘋果手機之人，遲早須對一靈魂之問：iCloud空間滿奈何？

\begin{itemize}
  \item \textbf{清理}：往設置 君之名 iCloud 管理儲存空間，視誰佔了君之空間。一般而言，照片與備份乃大頭。可先刪不需之舊備份。
  \item \textbf{升級}：50GB每月6塊錢，200GB每月21塊錢。個人覺50GB於多數人足矣，主要用以備份與存照片之縮略圖。
  \item \textbf{曲線救國}：若君不欲花錢，可將照片定期導出至電腦上，而後關iCloud照片同步。然如此，君即失照片自動備份之便。
\end{itemize}

一重要之提醒：iCloud照片圖庫非“網盤”——乃“同步”。君於手機上刪一張照片，iCloud上亦將刪。若君所欲乃“備份至雲端而後清空手機”，則iCloud照片非君所欲，君需乃iCloud雲備份，或第三方之網盤。

\begin{tcbinfo}{iCloud+：付費用戶才能享受之隱私三件套}
只要升級到任意付費iCloud檔位（50GB/200GB/2TB等），遂自動成爲\textbf{iCloud+}用戶，解鎖三個隱私護能：

\begin{itemize}
  \item \textbf{Private Relay（私密中繼）}：蘋果版之"VPN"。Safari瀏覽之流量經蘋果服務器中轉後再發往網站，沒有人（包括蘋果）能同時看到汝乃誰、汝訪問矣何。適合於公共Wi-Fi下保護隱私。設置路徑：\texttt{設置 > 汝之名字 > iCloud > 私密中繼}。
  \item \textbf{Hide My Email（隱郵件）}：於網站/App註冊時生成一次性隨機郵箱地址，郵件將自動轉發至汝之真實郵箱。不欲再收這家網站之郵件，直接關掉此隨機地址遂行，真實郵箱不暴露。設置路徑：\texttt{設置 > 汝之名字 > iCloud > 隱藏郵件}。
  \item \textbf{自定義電子郵件域名}：可用汝自己之域名（比如@汝之姓.com）收發iCloud郵件，適合個人品牌或家庭郵箱管理。
\end{itemize}

提示：Private Relay於某些國家/地區不可用，國內用戶或須其他方案補充。

\end{tcbinfo}

\begin{tcbwarn}{iCloud高級數據保護（ADP）——真正之端至端加密}
2022年底蘋果推出矣\textbf{iCloud高級數據護（Advanced Data Protection，ADP）}，以端至端加密（E2EE）自原來之健康、密碼等少數數據，擴展至了\textbf{絕大部iCloud數據}——照片、備忘錄、iCloud雲備份、設信息等共23類數據。

\begin{itemize}
  \item \textbf{何乃E2EE}：加密密鑰唯存在汝之設備上，連蘋果服務器皆沒有。就算蘋果爲黑客攻破，或執法部門要求調取，蘋果亦拿不出能解密之明文數據。
  \item \textbf{開啓之徑}：\texttt{設置 > 汝之名字 > iCloud > 高級數據保護}，依示設復聯繫人/復密鑰（\textbf{必作}，丟矣密鑰汝之數據亦救不回來）。
  \item \textbf{注意之事}：(1)開啓後務必妥善保管恢復密鑰，否則賬號丟失時無法找回數據；(2)共享相簿、iCloud共享圖庫等少數功能不參與E2EE；(3)中國大陸區Apple ID暫不支持ADP，須切換至其他地區賬號。
\end{itemize}

一句話建議：若汝之賬號裏有重要隱私數據，且賬號地區支持，強烈建議開啓ADP——這乃普通人能獲得之最強雲端數據保護。

\end{tcbwarn}

\booksection{四、同步與備份，勿混之}

此二詞聽來相仿，然則果全異：

\begin{itemize}
  \item \textbf{同步}乃雙向者：君於手機上刪之，雲端亦刪之；君於雲端刪之，手機上亦隨之刪。iCloud照片、通訊錄、備忘錄此等皆乃同步。利乃多設備之間保持一致，弊乃"刪即真無矣"。
  \item \textbf{備份}乃單向者：手機將當前狀態拍一"快照"存至雲端。後君於手機上刪物，雲端備份中之那份尚在。iCloud雲備份、電腦本地備份，皆乃備份。利乃數據更安全，弊乃佔空間，且不能實時更新。
\end{itemize}

\booksection{五、安卓用戶之賬號矩陣}

安卓這邊之情較蘋果複雜些許，因君之數據或散於數個賬號中：

\begin{itemize}
  \item \textbf{品牌賬號}（華爲/小米/OPPO/vivo）：管君之雲空間、查找手機、應用商店、錢包。每品牌皆有己之雲服務，通常乃基礎免費+擴容收費之式。
  \item \textbf{谷歌賬號}：若君用Gmail、Google相冊、Google通訊錄，此等物皆在谷歌賬號中。谷歌免費予15GB，較蘋果大方，然亦與Gmail、Drive共用。
  \item \textbf{微信/QQ}：其乃獨立於手機系統者。君之聊天記錄與微信藏，既不在品牌賬號中亦不在谷歌賬號中，只存於微信己之服務器上。
\end{itemize}

故安卓用戶作備份之時，須分別查：品牌雲備份、谷歌備份、微信聊天記錄備份，三者缺一不可。

\booksection{六、家庭共享：令家人亦受益}

蘋果之“家人共享”可使君與家人共享App購買、音樂、電影、iCloud空間，尚能共享位置。最多可加6人。設之之法：設置 君之名 家人共享。華爲與小米亦有類似之能，於賬號設置中能覓“家共享”或“家庭雲”。

一甚實用之場景：爲爸媽之手機上設好家人共享，如此君即可隨時見其位置（前提乃其同意），萬一老人走失或手機聯繫不上，能有一定位。

\begin{tcbwarn}{賬號被盜之後果遠較君所想爲重}
賬號被盜之後果遠較君所想爲重

\end{tcbwarn}

\begin{tcbpractice}{實操：爲君之賬號作一次全面體檢}
\begin{enumerate}
  \item 開手機設置，入賬號頁，確認雙重認證已啓。
  \item 開手機設置，入賬號頁，確認雙重認證已啓。
  \item 入iCloud或雲空間管理，視何數據在同步，關不需者。
  \item 確認"查找吾之手機"已啓。
  \item 確認"查找吾之手機"已啓。
\end{enumerate}

\end{tcbpractice}
